\documentclass[runningheads]{llncs}
%
\usepackage[T1]{fontenc}
\usepackage[utf8]{inputenc}
\usepackage{graphicx}

\begin{document}
%
\title{Título do artigo para o LNCS}
%
\titlerunning{Título curto}
%
\author{Primeiro Autor\inst{1}\orcidID{0000-0000-0000-0000} \and
Segundo Autor\inst{2}}
%
\authorrunning{P. Autor e S. Autor}
%
\institute{Instituição Um, Cidade, País\\
\email{autor1@exemplo.org} \and
Instituição Dois, Cidade, País\\
\email{autor2@exemplo.org}}
%
\maketitle
%
\begin{abstract}
Escreva aqui o resumo do artigo, com cerca de 150 a 250 palavras: o problema, o método,
os principais resultados e a conclusão.
\keywords{Primeira palavra-chave \and Segunda palavra-chave \and Terceira palavra-chave.}
\end{abstract}

\section{Introdução}
Este é um modelo mínimo para artigos no formato Lecture Notes in Computer Science (LNCS), da Springer.
Substitua o texto pelo seu. A classe \texttt{llncs} cuida da formatação, e as referências usam o estilo
\texttt{splncs04}~\cite{exemplo}.

\section{Método}
Descreva o método. Equações, figuras e tabelas seguem os comandos habituais do \LaTeX.

\begin{equation}
  E = mc^2
\end{equation}

\section{Resultados}
Apresente os resultados. Para citar outras seções, use \verb|\label| e \verb|\ref|, como em
a Seção~\ref{sec:conclusao}.

\section{Conclusão}\label{sec:conclusao}
Resuma as conclusões e os trabalhos futuros.

\bibliographystyle{splncs04}
\bibliography{referencias}
\end{document}
